\section{Limitations and Research Agenda}

Jankurai does not prove full semantic correctness. A merge witness can show that changed paths were routed to owners, required proof lanes ran, receipts exist, and missing evidence was handled; it cannot prove every business invariant or exclude malicious compiler, runtime, or organizational governance compromise. Current evidence is conformance-oriented and repository-local. The reference implementation is young, some detectors are advisory or partial, and signed receipts, transparent evidence logs, cross-repo federation, and organizational trust remain future work.

\subsection{Threats to Validity}

The taxonomy may reflect source-material bias. Detectors may have false positives and false negatives. Partial rows can be overinterpreted as stronger coverage than they deserve. Reference-profile scoring can bias readers toward Rust/TypeScript/PostgreSQL even when another stack can conform. Browser evidence is environment-sensitive. Screenshots, logs, and traces can leak secrets if redaction fails. Scores can be gamed. Compliance theater can produce artifacts without improving engineering. CI cost can become excessive. Small repositories may need a thinner profile than the paper emphasizes.

\subsection{Failure Modes of Jankurai Itself}

Jankurai can fail as artifact theater: reports exist but nobody trusts or uses them. Maps can go stale. Receipts can be faked or copied. Proof lanes can become flaky. Adapters can drift. Reference-profile guidance can harden into stack dogma. Ceremony can overload small changes. These are not reasons to abandon conformance; they are reasons to keep proof freshness, receipt validity, waiver expiry, seed fixtures, and independent implementation tests central.

Important research questions remain open. Which rules best predict escaped defects rather than reviewer discomfort? How much proof is enough for small changes? Which evidence classes can be public without leaking sensitive data? How should registries prevent badge gaming? Which repair queues actually reduce time-to-fix? Jankurai's mitigation strategy is to keep claims versioned, evidence local, waivers expiring, profile guidance non-normative, and independent implementations testable against fixtures.

\begin{table*}[t]
\caption{Concrete future work for Jankurai.}
\label{tab:future-work}
\centering
\scriptsize
\begin{tabularx}{\textwidth}{L{0.22\textwidth} Y}
\toprule
\JKTableHeader
\textbf{Lane} & \textbf{Expected artifact} \\
\midrule
Public conformance suite & Runnable tests, fixtures, expected JSON/Markdown reports, and versioned pass criteria. \\
Hosted or static evidence registry & Immutable evidence bundles, badges, score history, and current/previous standard metadata. \\
Signed proof receipts and attestations & Artifact digests, command receipts, provenance, and verification command. \\
Benchmark corpus for repair locality & Repos, patches, route plans, expected lanes, repair time, and escaped-defect labels. \\
Token-budget optimizer & Short root routers, path-scoped policy packs, symbol/repo maps, filtered command receipts, generated context packs, evidence indexes, stale-doc pruning, and deterministic proof routing. \\
Context-pack minimizer & Trust-labeled packs with measured token cost, recall, and false-exclusion rates. \\
Cross-agent adapter conformance tests & Codex, Claude, Cursor, Copilot, terminal-agent, hook, and MCP adapter fixtures. \\
UI geometry oracle hardening & Viewport matrices, screenshot crops, ARIA snapshots, layout rules, and false-positive corpus. \\
Semantic rule detector expansion & Fixtures and negative tests for authz, input boundaries, cost controls, release readiness, and human-review evidence. \\
Waiver and exception libraries & Runtime exception templates plus standards-waiver templates with docs links and repair hints. \\
Cost-budget and kill-switch tooling & Spend limits, quotas, stop conditions, alerts, and receipt-backed emergency controls. \\
Reusable cell registry and certification & Cell manifests, install plans, proof lanes, benchmark receipts, and certification policy. \\
\bottomrule
\end{tabularx}
\end{table*}

The strongest version of the criticism is useful: a standard can become ceremony, and a reference profile can become dogma. Jankurai should be judged by whether it improves repair locality, reduces unsafe merges, shortens review cycles, reduces vibe-artifact density, and produces better evidence. If a team can prove those outcomes with a different stack and a compatible control plane, the profile should be documented and studied.

\section{Conclusion}

Programming history improved expression, abstraction, runtime safety, and ecosystem reach. It did not solve repository intent. The hard questions remain social and technical at once: who owns this behavior, which boundary is authoritative, what proof is fresh, which waiver is still valid, and whether a refactor changed behavior.

AI makes ambiguity cheaper to create. It can write plausible code in the local dialect faster than a team can review by memory. That does not make agents the enemy; it makes agent-first repository design necessary. Ownership, generated boundaries, proof lanes, stable errors, repair receipts, and merge witnesses must be explicit enough that humans and agents work against the same evidence surface.

The durable contribution is not a stack ranking. It is continuous repository alignment: changed paths, owner route, required proof, observed evidence, missing-evidence decision, artifact digests, commit identity, tool identity, waiver state, and repair queue bound into a merge witness. Generation is cheap; reproducible merge evidence is the new standard. No proof, no merge. No receipt, no trust.
